\documentclass[10pt,aspectratio=43]{beamer}

\usepackage[english]{babel}
\usepackage{fontspec}
\setmainfont{Latin Modern Roman}
\setsansfont{Latin Modern Sans}
\setmonofont{Cascadia Code}[Scale=0.88]
\usepackage{amsmath, amssymb}
\usepackage{booktabs}
\usepackage{graphicx}
\usepackage{xcolor}
\usepackage{hyperref}
\usepackage{listings}

\usetheme{default}
\usecolortheme{default}
\setbeamertemplate{navigation symbols}{}
\setbeamertemplate{caption}[numbered]
\setbeamersize{text margin left=8mm,text margin right=8mm}

\definecolor{projectblue}{RGB}{24, 70, 125}
\definecolor{projectgray}{RGB}{82, 91, 103}

\setbeamercolor{structure}{fg=projectblue}
\setbeamercolor{frametitle}{fg=projectblue}
\setbeamercolor{title}{fg=projectblue}
\setbeamercolor{normal text}{fg=black}
\setbeamercolor{block title}{fg=white,bg=projectblue}
\setbeamercolor{block body}{fg=black,bg=projectblue!5}
\setbeamercolor{footline}{fg=projectgray}

\setbeamertemplate{frametitle}{%
    \vspace{0.45em}
    \insertframetitle\par
    \vspace{0.25em}
    {\color{projectblue!35}\rule{\linewidth}{0.5pt}}
}
\setbeamertemplate{footline}{%
    \hfill
    \usebeamercolor[fg]{footline}
    \scriptsize\insertframenumber/\inserttotalframenumber
    \hspace{0.8em}\vspace{0.45em}
}

\lstdefinestyle{cstyle}{
    language=C,
    basicstyle=\ttfamily\scriptsize,
    keywordstyle=\color{blue!60!black},
    commentstyle=\color{green!40!black},
    stringstyle=\color{red!50!black},
    showstringspaces=false,
    breaklines=true,
    columns=fullflexible,
    keepspaces=true,
    frame=single,
    rulecolor=\color{black!20},
    backgroundcolor=\color{black!2}
}

\lstdefinestyle{shellstyle}{
    language=bash,
    basicstyle=\ttfamily\scriptsize,
    keywordstyle=\color{blue!60!black},
    commentstyle=\color{green!40!black},
    stringstyle=\color{red!50!black},
    showstringspaces=false,
    breaklines=true,
    columns=fullflexible,
    keepspaces=true,
    frame=single,
    rulecolor=\color{black!20},
    backgroundcolor=\color{black!2}
}

\newcommand{\projecttitle}{heapx}
\newcommand{\projectsubtitle}{A Compact Experimental C Library for Heap Data Structures}

\title{\projecttitle}
\subtitle{\projectsubtitle}
\author{heapx project}
\institute{C99 data structures and algorithm engineering}
\date{\today}

\begin{document}

\begin{frame}{Title and Objective}
    \vspace{0.4em}
    {\Large\bfseries\color{projectblue}\projecttitle}

    \vspace{0.4em}
    {\large \projectsubtitle}

    \vspace{1.1em}
    \begin{block}{Objective}
        Compare multiple heap implementations behind one stable C API while making API contracts, correctness checks, memory behavior, and benchmarks explicit.
    \end{block}

    \vfill
    \small heapx project \hfill C99 data structures and algorithm engineering
\end{frame}

\begin{frame}{Motivation}
    Priority queues are central in graph algorithms, scheduling, simulations, and optimization.

    \vspace{0.5em}
    heapx studies the engineering side of heap implementations while staying intentionally narrow:
    \begin{itemize}
        \item one C99 heap-shaped API, not a broad container library;
        \item private backend representations selected at creation time;
        \item array storage versus pointer-heavy node structures;
        \item worst-case versus amortized complexity;
        \item targeted operations through handles;
        \item reproducible tests, invariants, documentation, and benchmarks.
    \end{itemize}

    \vspace{0.4em}
    \begin{block}{Guiding question}
        Can different heap backends share one compact API without hiding their trade-offs?
    \end{block}
\end{frame}

\begin{frame}[fragile]{Public API}
    The public API exposes an opaque heap handle and selects the backend only at creation time.

    \begin{lstlisting}[style=cstyle]
struct heapx_heap *heapx_create(
    enum heapx_implementation implementation,
    heapx_cmp_fn cmp
);

int heapx_insert(struct heapx_heap *heap, void *item);
int heapx_insert_handle(
    struct heapx_heap *heap,
    void *item,
    struct heapx_handle *out
);
int heapx_decrease_key(struct heapx_heap *heap,
                       struct heapx_handle handle);
void *heapx_remove(struct heapx_heap *heap,
                   struct heapx_handle handle);
    \end{lstlisting}
\end{frame}

\begin{frame}{API Contract}
    The API contract is deliberately explicit:

    \begin{itemize}
        \item stored objects are caller-owned \texttt{void *} pointers;
        \item ordering is defined by \texttt{heapx\_cmp\_fn};
        \item equal-priority extraction order is not stable;
        \item \texttt{NULL} heap behavior is defined in public functions;
        \item \texttt{contains} is linear and diagnostic, not heap-native performance API;
        \item no internal synchronization is provided.
    \end{itemize}

    \vspace{0.5em}
    The public layer owns common edge cases so backends can focus on algorithms.
\end{frame}

\begin{frame}{Complexity Summary}
    The implemented backends occupy different points in the heap design space:
    array-based worst-case bounds, classic amortized structures, and a simpler Fibonacci-style variant.

    \vspace{0.4em}
    \begin{center}
        \begin{tabular}{llll}
            \toprule
            Backend        & Insert      & Decrease key & Extract min     \\
            \midrule
            Binary heap    & $O(\log n)$ & $O(\log n)$  & $O(\log n)$     \\
            Fibonacci heap & am. $O(1)$  & am. $O(1)$   & am. $O(\log n)$ \\
            Kaplan heap    & am. $O(1)$  & am. $O(1)$   & am. $O(\log n)$ \\
            \bottomrule
        \end{tabular}
    \end{center}

    \vspace{0.6em}
    \begin{itemize}
        \item \texttt{remove}: $O(\log n)$ for binary, amortized $O(\log n)$ for Fibonacci-style heaps.
        \item \texttt{peek\_min}, \texttt{size}, \texttt{empty}: $O(1)$.
        \item \texttt{contains}: $O(n)$ in all current backends.
    \end{itemize}
\end{frame}

\begin{frame}{Internal Architecture}
    \begin{block}{Core idea}
        One public dispatch layer, multiple private backend implementations.
    \end{block}

    \vspace{0.4em}
    \begin{itemize}
        \item Concrete heaps embed \texttt{struct heapx\_heap} as their first field.
        \item The base object stores comparator, vtable, heap id, and handle slots.
        \item Public functions validate common behavior and dispatch through the vtable.
        \item Backends own their storage, repair algorithms, and invariant checks.
    \end{itemize}
\end{frame}

\begin{frame}[fragile]{Vtable Model}
    The internal vtable mirrors the public heap operations.

    \begin{lstlisting}[style=cstyle]
struct heapx_vtable {
    void (*destroy)(struct heapx_heap *heap);
    int (*insert)(struct heapx_heap *heap, void *item);
    int (*decrease_key)(struct heapx_heap *heap,
                        void *owner);
    void *(*remove)(struct heapx_heap *heap,
                    void *owner);
    void *(*extract_min)(struct heapx_heap *heap);
    int (*check_invariants)(const struct heapx_heap *heap);
};
    \end{lstlisting}

    Targeted operations receive an already resolved backend owner pointer.
\end{frame}

\begin{frame}{Generational Handles}
    Decrease-key and arbitrary remove need a stable reference to an item position.

    \vspace{0.4em}
    A heapx handle contains:
    \begin{itemize}
        \item logical heap id;
        \item slot index in the common handle table;
        \item slot generation.
    \end{itemize}

    \vspace{0.5em}
    \begin{block}{Safety property}
        Stale handles and handles from another heap fail cleanly, even when internal slots are reused.
    \end{block}
\end{frame}

\begin{frame}{Handle Owners}
    The common handle table maps a public handle to a backend owner pointer.

    \begin{center}
        \begin{tabular}{lll}
            \toprule
            Backend        & Owner pointer  & Why                \\
            \midrule
            Binary heap    & Stable locator & Array entries move \\
            Fibonacci heap & Node pointer   & Nodes are stable   \\
            Kaplan heap    & Node pointer   & Nodes are stable   \\
            \bottomrule
        \end{tabular}
    \end{center}

    \vspace{0.5em}
    This keeps public handle validation common while allowing backend-specific repair logic.
\end{frame}

\begin{frame}{Binary Heap Backend}
    Binary heap implementation:

    \begin{itemize}
        \item dynamic array of entries;
        \item zero-based parent/child index arithmetic;
        \item \texttt{sift\_up} after insertion and decrease-key;
        \item \texttt{sift\_down} after extract-min;
        \item \texttt{repair\_at} after arbitrary remove;
        \item stable locators for handled entries.
    \end{itemize}

    \vspace{0.5em}
    \begin{block}{Trade-off}
        Weaker decrease-key bound than Fibonacci-style heaps, but simple code and cache-friendly storage.
    \end{block}
\end{frame}

\begin{frame}{Fibonacci Heap Backend}
    Fibonacci heap implementation:

    \begin{itemize}
        \item circular root and child lists;
        \item parent pointers, degree, and mark bits;
        \item insert adds a singleton root;
        \item decrease-key performs cut plus cascading cuts;
        \item extract-min promotes children and consolidates roots by degree;
        \item reusable degree table for consolidation.
    \end{itemize}

    \vspace{0.5em}
    Designed around amortized $O(1)$ insert and decrease-key.
\end{frame}

\begin{frame}{Kaplan Heap Backend}
    Kaplan heap implementation:

    \begin{itemize}
        \item simple Fibonacci heap variant;
        \item steady state is one heap-ordered tree;
        \item insertion links a singleton node against the root;
        \item decrease-key performs cascading rank/state changes and one cut;
        \item extract-min consolidates children with fair links by rank;
        \item remove structurally simulates decrease-to-minus-infinity plus extract-min.
    \end{itemize}

    \vspace{0.5em}
    It is algorithmically close to Fibonacci heaps, but the repair mechanism is different.
\end{frame}

\begin{frame}{Memory and Allocation Safety}
    heapx separates caller-owned items from heap-owned metadata.

    \vspace{0.4em}
    \begin{itemize}
        \item Stored items are never copied or freed by heapx.
        \item Allocation size arithmetic uses centralized overflow checks.
        \item Pointer-heavy backends use fixed-size node pools.
        \item Pool slots are aligned for backend node structs.
        \item Temporary consolidation allocation failures preserve correctness through fallback behavior.
    \end{itemize}

    \vspace{0.4em}
    The goal is predictable failure behavior, not only fast successful paths.
\end{frame}

\begin{frame}{Correctness Strategy}
    The test strategy combines local and global checks.

    \begin{itemize}
        \item deterministic tests for edge cases and public API behavior;
        \item randomized oracle test with inserts, decrease-key, remove, extract and peek;
        \item common invariant checks for handles and base state;
        \item backend-specific invariant checks;
        \item ASan, UBSan, leak-check and optimized builds;
        \item Dijkstra checksum agreement across backends.
    \end{itemize}
\end{frame}

\begin{frame}{Invariant Examples}
    \begin{itemize}
        \item \textbf{Common layer}: valid vtable, comparator, heap id, handle slot free list.
        \item \textbf{Binary heap}: parent-child order, locator index consistency, handle resolution.
        \item \textbf{Fibonacci heap}: circular links, parent pointers, degree counts, minimum root.
        \item \textbf{Kaplan heap}: single root, rank constraints, sibling links, cleared rank table.
    \end{itemize}

    \vspace{0.6em}
    Invariants are run after successful mutations when internal checks are enabled.
\end{frame}

\begin{frame}[fragile]{Benchmark Methodology}
    Two benchmark layers:

    \begin{itemize}
        \item \textbf{Heap-only}: isolates operation mixes such as insert-only, decrease-only, drain, and mixed handles.
        \item \textbf{Dijkstra}: exercises insert-handle, decrease-key and extract-min inside a shortest-path workload.
    \end{itemize}

    \vspace{0.4em}
    \begin{lstlisting}[style=shellstyle]
make benchmark-heap-tsv N=500000 SEED=123
make benchmark-heap-compare BASE=before.tsv HEAD=after.tsv

make benchmark GRAPH=graphs/dimacs/USA-road-d.USA.gr SOURCE=1
    \end{lstlisting}
\end{frame}

\begin{frame}{Benchmark Interpretation}
    Benchmarks are empirical evidence for this implementation and workload.

    \vspace{0.4em}
    Results depend on:
    \begin{itemize}
        \item compiler and optimization flags;
        \item machine and cache behavior;
        \item allocator behavior;
        \item graph structure and operation mix;
        \item frequency of decrease-key.
    \end{itemize}

    \vspace{0.4em}
    Smoke tests validate execution and consistency; large datasets are needed for strong performance claims.
\end{frame}

\begin{frame}{Current Status}
    \begin{block}{What is already in place}
        A compact public API over three heap backends, with tests, invariant checks, sanitizer targets, benchmark tooling, Doxygen documentation, report, and presentation.
    \end{block}

    \vspace{0.5em}
    \begin{itemize}
        \item The architecture is stable and backend-oriented.
        \item Targeted operations are safe through generational handles.
        \item Benchmarks can be saved and compared as TSV.
        \item The repository remains focused on heap data structures.
    \end{itemize}
\end{frame}

\begin{frame}{Limitations and Next Steps}
    The current state is professional enough to study, but still experimental.

    \vspace{0.4em}
    \begin{columns}[T,onlytextwidth]
        \begin{column}{0.48\textwidth}
            \textbf{Limitations}
            \begin{itemize}
                \item preliminary benchmark conclusions;
                \item linear \texttt{contains};
                \item no synchronization or custom allocator API;
                \item no stable ordering for equal priorities.
            \end{itemize}
        \end{column}
        \begin{column}{0.48\textwidth}
            \textbf{Next steps}
            \begin{itemize}
                \item larger fixed benchmark campaigns;
                \item reproducible benchmark summaries;
                \item optional analyzer CI jobs;
                \item continued invariant documentation.
            \end{itemize}
        \end{column}
    \end{columns}
\end{frame}

\begin{frame}{Takeaway}
    \begin{block}{Main takeaway}
        heapx is an algorithm-engineering workspace: one heap-shaped API, multiple private backends, and tooling to compare correctness and performance without changing client code.
    \end{block}

    \vspace{0.7em}
    For a 30-minute discussion, the interesting point is not only which heap is faster.
    It is how API design, handles, memory management, invariants, and benchmarks make that comparison trustworthy.
\end{frame}

\end{document}
